\documentclass[11pt]{article}

% macros.tex — Vegas paper notation

% ── Packages ──────────────────────────────────────────────────────────
\usepackage{amsmath,amssymb,amsthm}
\usepackage{stmaryrd}       % \llbracket, \rrbracket
\usepackage{mathtools}       % \coloneqq
\usepackage{xcolor}
\usepackage{hyperref}
\usepackage{cleveref}
\usepackage{booktabs}
\usepackage{enumitem}
\usepackage{mathpartir}      % pretty inference rules (operational semantics)
\usepackage{simplebnf}       % pretty BNF grammars

% ── Theorem environments ──────────────────────────────────────────────
\newtheorem{theorem}{Theorem}[section]
\newtheorem{lemma}[theorem]{Lemma}
\newtheorem{proposition}[theorem]{Proposition}
\newtheorem{corollary}[theorem]{Corollary}
\theoremstyle{definition}
\newtheorem{definition}[theorem]{Definition}
\newtheorem{example}[theorem]{Example}
\theoremstyle{remark}
\newtheorem{remark}[theorem]{Remark}
\newtheorem{notation}[theorem]{Notation}

% ── Players and sets ──────────────────────────────────────────────────
\newcommand{\Players}{\mathcal{I}}

% ── Types and values ──────────────────────────────────────────────────
\newcommand{\Ty}{\mathsf{Ty}}
\newcommand{\Val}{\mathsf{Val}}

% ── Visibility annotations ────────────────────────────────────────────
\newcommand{\pub}{\ensuremath{\mathsf{public}}}
\newcommand{\sealed}{\ensuremath{\mathsf{sealed}}}

% ── Contexts ──────────────────────────────────────────────────────────
\newcommand{\view}[1]{\Gamma\!\!\downarrow_{#1}}       % viewVCtx
\newcommand{\erasectx}[1]{\| #1 \|}                    % eraseVCtx
\newcommand{\erasepub}[1]{\| #1 \|^{\pub}}             % erasePubVCtx (public-only erasure operator)

% ── Constructs ────────────────────────────────────────────────────────
\newcommand{\Ret}{\ensuremath{\mathbf{ret}}}
\newcommand{\Let}{\ensuremath{\mathbf{let}}}
\newcommand{\Sample}{\ensuremath{\mathbf{sample}}}
\newcommand{\Commit}{\ensuremath{\mathbf{commit}}}
\newcommand{\Reveal}{\ensuremath{\mathbf{reveal}}}
\newcommand{\Where}{\ensuremath{\mathbf{where}}}

% ── Small-step operational semantics ──────────────────────────────────
\newcommand{\sconf}[3]{\langle #1 \mathrel{;} #2 \mathrel{;} #3 \rangle} % configuration
\newcommand{\sstep}{\longrightarrow}                  % step relation
\newcommand{\lstep}[1]{\xrightarrow{\,#1\,}}          % labelled step
\newcommand{\sem}[1]{\llbracket #1 \rrbracket}        % evaluation in L
\newcommand{\supp}{\mathop{\mathrm{supp}}}            % support of a weighting
\newcommand{\outcomeof}{\mathsf{outcome}}             % terminal outcome reader

% ── Distributions ─────────────────────────────────────────────────────
\newcommand{\FWeight}[1]{\mathcal{W}(#1)}
\newcommand{\PMF}[1]{\mathrm{PMF}(#1)}

% ── Outcomes ──────────────────────────────────────────────────────────
\newcommand{\Outcome}{\mathcal{O}}

% ── Profiles and strategies ───────────────────────────────────────────
\newcommand{\Profile}{\sigma}

% ── Strategic bridge ──────────────────────────────────────────────────
\newcommand{\KG}{\ensuremath{\mathcal{G}}}
\newcommand{\EU}{\ensuremath{\mathrm{EU}}}

% ── Environments ──────────────────────────────────────────────────────
\newcommand{\Env}{\ensuremath{\mathrm{Env}}}

% ── Misc math ─────────────────────────────────────────────────────────
\newcommand{\NNQ}{\ensuremath{\mathbb{Q}_{\ge 0}}}
\newcommand{\bind}{\mathbin{\gg\!=}}

% ── Lean references ───────────────────────────────────────────────────
\newcommand{\TT}[1]{\texttt{#1}}

% ── Event graphs and configurations ───────────────────────────────────
\newcommand{\EG}{\mathcal{E}}                       % an event graph
\newcommand{\prereq}{\mathrm{prereq}}
\newcommand{\Cfg}{\mathsf{Cfg}}                     % configurations of an EG
\newcommand{\readsOf}{\mathrm{reads}}

% ── Machines ──────────────────────────────────────────────────────────
\newcommand{\State}{\mathsf{S}}
\newcommand{\Action}{\mathsf{A}}

% ── Solution concepts / kernels ───────────────────────────────────────
\newcommand{\Kbeh}{K^{\mathrm{beh}}}
\newcommand{\Kpure}{K^{\mathrm{pure}}}
\newcommand{\Kmix}{K^{\mathrm{mix}}}

% ── Faithfulness-relation shapes (the "forest" vocabulary) ────────────
% Symbols used in the guarantees map: an equivalence, a (one-directional)
% adequacy, and a refinement order.
\newcommand{\eqrel}{\cong}            % equivalence (symmetric)
\newcommand{\adeq}{\vDash}            % adequacy: right object IS the denotation
\newcommand{\ordrel}{\sqsubseteq}     % refinement order

% ── Non-intrusive "status / not-yet-mechanized" note ──────────────────
% A footnote that flags an aspirational or partial item without breaking
% the surrounding narrative.  Used sparingly.
\newcommand{\statusnote}[1]{\footnote{\textit{Status.}~#1}}

% ── Boxes for design notes / Lean correspondence ──────────────────────
\newcommand{\leancorr}[1]{%
  \par\smallskip
  {\raggedright\hbadness=10000\emergencystretch=3em
   \noindent\emph{Lean correspondence.} #1\par}%
  \smallskip}

% ── Lean signature boxes ──────────────────────────────────────────────
% A light-shaded, thin-ruled block for displaying Lean signatures.
% Contents are typeset (not verbatim) so we can mix \TT, math, and Greek
% as the rest of the document already does.
\usepackage[most]{tcolorbox}
\tcbset{
  leanstyle/.style={
    colback=black!3, colframe=black!30, boxrule=0.3pt, arc=1pt,
    left=8pt, right=8pt, top=4pt, bottom=4pt,
    breakable, enhanced jigsaw,
    before skip=4pt, after skip=4pt,
    fontupper=\footnotesize\ttfamily,
  },
}
\newtcolorbox{leanbox}[1][]{leanstyle, #1}


\usepackage[margin=1in]{geometry}

\title{VegasCore: Event-Graph-Native Protocol Semantics}
\author{}
\date{}

\begin{document}
\maketitle

\begin{abstract}
VegasCore is a Lean~4 mechanization of finite games written as programs.
Checked Vegas syntax is elaborated to an invariant-carrying event graph; that
graph is interpreted as a probabilistic, observation-aware machine; sequential
FOSG views and kernel games are derived from the machine. The supported
semantic path is
\[
  \WFP \longrightarrow \leanref{EventGraph}
  \longrightarrow \leanref{Machine}
  \longrightarrow \leanref{Machine.FOSGView}
  \longrightarrow \leanref{KernelGame}.
\]
The source language is the user-facing way to define games. The event-graph
machine is the semantic owner.
\end{abstract}

\tableofcontents

\section{Checked Programs}

The public boundary is \leanref{WFProgram}. A checked program contains the
root context, source program, initial environment, and facts for
freshness, scoped visibility, reveal completeness, sample normalization, and
nonempty legal commit guards.

Visibility is represented in the program context. Bindings are public or
hidden and owned by a player. Commit strategies and observations are indexed by
the information available to the acting player, so hidden information owned by
other players is not available to legal strategies.

\section{Protocol Denotation}

The executable protocol stack is concentrated in \leanref{Vegas.Protocol}.

\begin{center}
\small
\begin{tabular}{@{}ll@{}}
\toprule
File & Role \\
\midrule
\leanref{Protocol/EventGraph.lean} &
finite typed event graph \\
\leanref{Protocol/Machine.lean} &
probabilistic observation-aware machine carrier \\
\leanref{Protocol/EventGraphMachine.lean} &
event-graph configurations as machine states \\
\leanref{Protocol/ProgramEventGraph.lean} &
checked syntax elaborated to an event graph, with finiteness and observability facts \\
\leanref{Protocol/Trace.lean} &
bounded event/state traces and machine outcome kernels \\
\leanref{Protocol/FOSG.lean} &
sequential FOSG views derived from machines \\
\leanref{Protocol/Strategic.lean} &
finite program event-graph kernel-game constructors \\
\leanref{Protocol/Kuhn.lean} &
machine-native and Vegas-facing Kuhn realization theorems \\
\leanref{Protocol/Backend.lean} &
machine-to-machine backend refinement obligations \\
\bottomrule
\end{tabular}
\end{center}

An event graph configuration is an extensional partial assignment of graph
node results, together with the invariant that completed nodes are closed under
dependencies. The frontier is computed from the assignment; it is not cached as
machine state. A machine step executes one enabled event: a player-owned
frontier node via a legal write slice, or a non-player frontier node via the
internal event channel.

\begin{theorem}[Frontier stability]
Let \(C\) be an event graph configuration. If \(a\) and \(b\) are distinct
nodes on the frontier of \(C\), then after executing \(a\) with any legal result
slice, \(b\) is still on the frontier of the successor configuration.

In Lean this is \leanref{EventGraph.Configuration.withResult_mem_frontier_of_ne}.
\end{theorem}

\begin{theorem}[Frontier diamond]
Let \(a\) and \(b\) be distinct frontier nodes of a configuration \(C\). Let
\(s_a\) and \(s_b\) be legal result slices for \(a\) and \(b\), respectively.
Then the two linearizations reach the same extensional configuration:
\[
  (C[a \mapsto s_a])[b \mapsto s_b]
  =
  (C[b \mapsto s_b])[a \mapsto s_a].
\]

In Lean this is \leanref{EventGraph.Configuration.withResult_comm}.
\end{theorem}

Thus a user may read a frontier execution linearly without losing semantic
content: the linear order is presentation data. Intermediate transcripts may
record which enabled event happened first, but independent frontier events cannot
race to change the final physical state. Any strategically relevant dependency
must already be represented as an event-graph dependency, and therefore prevents the
dependent event from being on the same frontier.

\section{FOSG As A View}

\leanref{Machine.FOSGView} derives a game-theoretic checkpoint view from a
protocol machine. Worlds are machine states; bounded views add only a depth
counter. A FOSG transition is a player-facing frontier round with its own
round-action alphabet, while the underlying machine remains the executable
primitive-event carrier. Thus FOSG is the strategic interface to the same
machine state space, not a second execution semantics.

The finite checked-program view is obtained from
\leanref{eventGraphFOSGView}. Its bounded version at \leanref{syntaxSteps}
provides the information-state-indexed strategy carriers used by the public
kernel games.

\section{Strategic API}

For a finite checked program \(g\), the public strategic carriers are:
\[
\begin{array}{ll}
\leanref{pureStrategyAt}\;g\;i &
\text{pure reachable-legal strategy for player } i,\\
\leanref{behavioralStrategyPMFAt}\;g\;i &
\text{PMF behavioral reachable-legal strategy for player } i,\\
\leanref{pureKernelGame}\;g &
\text{pure strategic-form kernel game},\\
\leanref{pmfBehavioralKernelGame}\;g &
\text{PMF behavioral strategic-form kernel game}.
\end{array}
\]

Equilibrium, dominance, welfare, potential-game, zero-sum, and related
predicates in \leanref{Vegas.Equilibrium}, \leanref{Vegas.PureStrategic}, and
\leanref{Vegas.GameProperties} are wrappers over these kernel games. Illegal
deviations are excluded by the strategy carrier.

\section{Kuhn Realization}

The main Vegas-facing finite mixed-to-behavioral theorem is
\leanref{kuhn_finite}. In Lean shape:
\[
\begin{array}{l}
\forall \mu : \prod_i
  \PMF{(\leanref{pureKernelGameAt}\;g).\leanref{Strategy}\;i},\\
\quad \exists \beta : (\leanref{pmfBehavioralKernelGameAt}\;g).\leanref{Profile},\\
\quad
  (\leanref{pmfBehavioralKernelGameAt}\;g).\leanref{outcomeKernel}\;\beta
  =
  \leanref{Math.PMFProduct.pmfPi}\;\mu
    \bind
    (\lambda \pi,\,
      (\leanref{pureKernelGameAt}\;g).\leanref{outcomeKernel}\;\pi).
\end{array}
\]

The corresponding named protocol property is \leanref{KuhnPMF}, realized by
\leanref{kuhnPMF_finite}. These are stated in Vegas kernel-game terms. The
bounded FOSG view is the finite sequential game to which Kuhn's theorem
applies.

\section{Public Interfaces}

The user-facing flow is:

\begin{enumerate}[nosep]
\item users write checked Vegas programs;
\item checked syntax elaborates to \leanref{EventGraph};
\item the graph is interpreted as \leanref{Machine};
\item FOSG views derive finite sequential-game structure from the machine;
\item public kernel games and theorem statements use Vegas names.
\end{enumerate}

\end{document}
